\IfFileExists{stacks-project.cls}{%
\documentclass{stacks-project}
}{%
\documentclass{amsart}
}

% For dealing with references we use the comment environment
\usepackage{verbatim}
\newenvironment{reference}{\comment}{\endcomment}
%\newenvironment{reference}{}{}
\newenvironment{slogan}{\comment}{\endcomment}
\newenvironment{history}{\comment}{\endcomment}

% For commutative diagrams we use Xy-pic
\usepackage[all]{xy}

% We use 2cell for 2-commutative diagrams.
\xyoption{2cell}
\UseAllTwocells

% We use multicol for the list of chapters between chapters
\usepackage{multicol}

% This is generally recommended for better output
\usepackage{lmodern}
\usepackage[T1]{fontenc}

% For cross-file-references
\usepackage{xr-hyper}

% Package for hypertext links:
\usepackage{hyperref}

% For any local file, say "hello.tex" you want to link to please
% use \externaldocument[hello-]{hello}
\externaldocument[introduction-]{introduction}
\externaldocument[conventions-]{conventions}
\externaldocument[sets-]{sets}
\externaldocument[categories-]{categories}
\externaldocument[topology-]{topology}
\externaldocument[sheaves-]{sheaves}
\externaldocument[sites-]{sites}
\externaldocument[stacks-]{stacks}
\externaldocument[fields-]{fields}
\externaldocument[algebra-]{algebra}
\externaldocument[brauer-]{brauer}
\externaldocument[homology-]{homology}
\externaldocument[derived-]{derived}
\externaldocument[simplicial-]{simplicial}
\externaldocument[more-algebra-]{more-algebra}
\externaldocument[smoothing-]{smoothing}
\externaldocument[modules-]{modules}
\externaldocument[sites-modules-]{sites-modules}
\externaldocument[injectives-]{injectives}
\externaldocument[cohomology-]{cohomology}
\externaldocument[sites-cohomology-]{sites-cohomology}
\externaldocument[dga-]{dga}
\externaldocument[dpa-]{dpa}
\externaldocument[sdga-]{sdga}
\externaldocument[hypercovering-]{hypercovering}
\externaldocument[schemes-]{schemes}
\externaldocument[constructions-]{constructions}
\externaldocument[properties-]{properties}
\externaldocument[morphisms-]{morphisms}
\externaldocument[coherent-]{coherent}
\externaldocument[divisors-]{divisors}
\externaldocument[limits-]{limits}
\externaldocument[varieties-]{varieties}
\externaldocument[topologies-]{topologies}
\externaldocument[descent-]{descent}
\externaldocument[perfect-]{perfect}
\externaldocument[more-morphisms-]{more-morphisms}
\externaldocument[flat-]{flat}
\externaldocument[groupoids-]{groupoids}
\externaldocument[more-groupoids-]{more-groupoids}
\externaldocument[etale-]{etale}
\externaldocument[chow-]{chow}
\externaldocument[intersection-]{intersection}
\externaldocument[pic-]{pic}
\externaldocument[weil-]{weil}
\externaldocument[adequate-]{adequate}
\externaldocument[dualizing-]{dualizing}
\externaldocument[duality-]{duality}
\externaldocument[discriminant-]{discriminant}
\externaldocument[derham-]{derham}
\externaldocument[local-cohomology-]{local-cohomology}
\externaldocument[algebraization-]{algebraization}
\externaldocument[curves-]{curves}
\externaldocument[resolve-]{resolve}
\externaldocument[models-]{models}
\externaldocument[functors-]{functors}
\externaldocument[equiv-]{equiv}
\externaldocument[pione-]{pione}
\externaldocument[etale-cohomology-]{etale-cohomology}
\externaldocument[proetale-]{proetale}
\externaldocument[relative-cycles-]{relative-cycles}
\externaldocument[more-etale-]{more-etale}
\externaldocument[trace-]{trace}
\externaldocument[crystalline-]{crystalline}
\externaldocument[spaces-]{spaces}
\externaldocument[spaces-properties-]{spaces-properties}
\externaldocument[spaces-morphisms-]{spaces-morphisms}
\externaldocument[decent-spaces-]{decent-spaces}
\externaldocument[spaces-cohomology-]{spaces-cohomology}
\externaldocument[spaces-limits-]{spaces-limits}
\externaldocument[spaces-divisors-]{spaces-divisors}
\externaldocument[spaces-over-fields-]{spaces-over-fields}
\externaldocument[spaces-topologies-]{spaces-topologies}
\externaldocument[spaces-descent-]{spaces-descent}
\externaldocument[spaces-perfect-]{spaces-perfect}
\externaldocument[spaces-more-morphisms-]{spaces-more-morphisms}
\externaldocument[spaces-flat-]{spaces-flat}
\externaldocument[spaces-groupoids-]{spaces-groupoids}
\externaldocument[spaces-more-groupoids-]{spaces-more-groupoids}
\externaldocument[bootstrap-]{bootstrap}
\externaldocument[spaces-pushouts-]{spaces-pushouts}
\externaldocument[spaces-chow-]{spaces-chow}
\externaldocument[groupoids-quotients-]{groupoids-quotients}
\externaldocument[spaces-more-cohomology-]{spaces-more-cohomology}
\externaldocument[spaces-simplicial-]{spaces-simplicial}
\externaldocument[spaces-duality-]{spaces-duality}
\externaldocument[formal-spaces-]{formal-spaces}
\externaldocument[restricted-]{restricted}
\externaldocument[spaces-resolve-]{spaces-resolve}
\externaldocument[formal-defos-]{formal-defos}
\externaldocument[defos-]{defos}
\externaldocument[cotangent-]{cotangent}
\externaldocument[examples-defos-]{examples-defos}
\externaldocument[algebraic-]{algebraic}
\externaldocument[examples-stacks-]{examples-stacks}
\externaldocument[stacks-sheaves-]{stacks-sheaves}
\externaldocument[criteria-]{criteria}
\externaldocument[artin-]{artin}
\externaldocument[quot-]{quot}
\externaldocument[stacks-properties-]{stacks-properties}
\externaldocument[stacks-morphisms-]{stacks-morphisms}
\externaldocument[stacks-limits-]{stacks-limits}
\externaldocument[stacks-cohomology-]{stacks-cohomology}
\externaldocument[stacks-perfect-]{stacks-perfect}
\externaldocument[stacks-introduction-]{stacks-introduction}
\externaldocument[stacks-more-morphisms-]{stacks-more-morphisms}
\externaldocument[stacks-geometry-]{stacks-geometry}
\externaldocument[moduli-]{moduli}
\externaldocument[moduli-curves-]{moduli-curves}
\externaldocument[examples-]{examples}
\externaldocument[exercises-]{exercises}
\externaldocument[guide-]{guide}
\externaldocument[desirables-]{desirables}
\externaldocument[coding-]{coding}
\externaldocument[obsolete-]{obsolete}
\externaldocument[fdl-]{fdl}
\externaldocument[index-]{index}

% Theorem environments.
%
\theoremstyle{plain}
\newtheorem{theorem}[subsection]{Theorem}
\newtheorem{proposition}[subsection]{Proposition}
\newtheorem{lemma}[subsection]{Lemma}

\theoremstyle{definition}
\newtheorem{definition}[subsection]{Definition}
\newtheorem{example}[subsection]{Example}
\newtheorem{exercise}[subsection]{Exercise}
\newtheorem{situation}[subsection]{Situation}

\theoremstyle{remark}
\newtheorem{remark}[subsection]{Remark}
\newtheorem{remarks}[subsection]{Remarks}

\numberwithin{equation}{subsection}

% Macros
%
\def\lim{\mathop{\mathrm{lim}}\nolimits}
\def\colim{\mathop{\mathrm{colim}}\nolimits}
\def\Spec{\mathop{\mathrm{Spec}}}
\def\Hom{\mathop{\mathrm{Hom}}\nolimits}
\def\Ext{\mathop{\mathrm{Ext}}\nolimits}
\def\SheafHom{\mathop{\mathcal{H}\!\mathit{om}}\nolimits}
\def\SheafExt{\mathop{\mathcal{E}\!\mathit{xt}}\nolimits}
\def\Sch{\mathit{Sch}}
\def\Mor{\mathop{\mathrm{Mor}}\nolimits}
\def\Ob{\mathop{\mathrm{Ob}}\nolimits}
\def\Sh{\mathop{\mathit{Sh}}\nolimits}
\def\NL{\mathop{N\!L}\nolimits}
\def\CH{\mathop{\mathrm{CH}}\nolimits}
\def\proetale{{pro\text{-}\acute{e}tale}}
\def\etale{{\acute{e}tale}}
\def\QCoh{\mathit{QCoh}}
\def\Ker{\mathop{\mathrm{Ker}}}
\def\Im{\mathop{\mathrm{Im}}}
\def\Coker{\mathop{\mathrm{Coker}}}
\def\Coim{\mathop{\mathrm{Coim}}}

% Boxtimes
%
\DeclareMathSymbol{\boxtimes}{\mathbin}{AMSa}{"02}

%
% Macros for moduli stacks/spaces
%
\def\QCohstack{\mathcal{QC}\!\mathit{oh}}
\def\Cohstack{\mathcal{C}\!\mathit{oh}}
\def\Spacesstack{\mathcal{S}\!\mathit{paces}}
\def\Quotfunctor{\mathrm{Quot}}
\def\Hilbfunctor{\mathrm{Hilb}}
\def\Curvesstack{\mathcal{C}\!\mathit{urves}}
\def\Polarizedstack{\mathcal{P}\!\mathit{olarized}}
\def\Complexesstack{\mathcal{C}\!\mathit{omplexes}}
% \Pic is the operator that assigns to X its picard group, usage \Pic(X)
% \Picardstack_{X/B} denotes the Picard stack of X over B
% \Picardfunctor_{X/B} denotes the Picard functor of X over B
\def\Pic{\mathop{\mathrm{Pic}}\nolimits}
\def\Picardstack{\mathcal{P}\!\mathit{ic}}
\def\Picardfunctor{\mathrm{Pic}}
\def\Deformationcategory{\mathcal{D}\!\mathit{ef}}
\usepackage{cleveref}
\crefname{originallemma}{original lemma}{original lemmas}
\Crefname{originallemma}{Original lemma}{Original lemmas}
\makeatletter
\newlabel{lemma-R-finite-above-x}{{03JU}{source}{}{sourcepointer-1}{}}
\newlabel{lemma-R-finite-above-x@cref}{{[originallemma][0][]03JU}{[1][0][]source}{}{}{}}
\makeatother
\crefname{originallemma}{original lemma}{original lemmas}
\Crefname{originallemma}{Original lemma}{Original lemmas}
\makeatletter
\newlabel{lemma-U-finite-above-x}{{03JS}{source}{}{sourcepointer-2}{}}
\newlabel{lemma-U-finite-above-x@cref}{{[originallemma][0][]03JS}{[1][0][]source}{}{}{}}
\makeatother
\crefname{originallemma}{original lemma}{original lemmas}
\Crefname{originallemma}{Original lemma}{Original lemmas}
\makeatletter
\newlabel{lemma-weak-UR-finite-above-x}{{040U}{source}{}{sourcepointer-3}{}}
\newlabel{lemma-weak-UR-finite-above-x@cref}{{[originallemma][0][]040U}{[1][0][]source}{}{}{}}
\makeatother
\crefname{originalremark}{original remark}{original remarks}
\Crefname{originalremark}{Original remark}{Original remarks}
\makeatletter
\newlabel{remark-recall}{{03II}{source}{}{sourcepointer-4}{}}
\newlabel{remark-recall@cref}{{[originalremark][0][]03II}{[1][0][]source}{}{}{}}
\makeatother
\crefname{originalreference}{original reference}{original references}
\Crefname{originalreference}{Original reference}{Original references}
\makeatletter
\newlabel{schemes-lemma-mono-towards-spec-field}{{03DP}{source}{}{sourcepointer-5}{}}
\newlabel{schemes-lemma-mono-towards-spec-field@cref}{{[originalreference][0][]03DP}{[1][0][]source}{}{}{}}
\makeatother
\crefname{originalreference}{original reference}{original references}
\Crefname{originalreference}{Original reference}{Original references}
\makeatletter
\newlabel{spaces-morphisms-lemma-composition-quasi-compact}{{03HG}{source}{}{sourcepointer-6}{}}
\newlabel{spaces-morphisms-lemma-composition-quasi-compact@cref}{{[originalreference][0][]03HG}{[1][0][]source}{}{}{}}
\makeatother
\crefname{originalreference}{original reference}{original references}
\Crefname{originalreference}{Original reference}{Original references}
\makeatletter
\newlabel{spaces-morphisms-lemma-quasi-compact-local}{{03KG}{source}{}{sourcepointer-7}{}}
\newlabel{spaces-morphisms-lemma-quasi-compact-local@cref}{{[originalreference][0][]03KG}{[1][0][]source}{}{}{}}
\makeatother
\crefname{originalreference}{original reference}{original references}
\Crefname{originalreference}{Original reference}{Original references}
\makeatletter
\newlabel{spaces-properties-lemma-equivalence-class-point-monomorphism}{{0H2X}{source}{}{sourcepointer-8}{}}
\newlabel{spaces-properties-lemma-equivalence-class-point-monomorphism@cref}{{[originalreference][0][]0H2X}{[1][0][]source}{}{}{}}
\makeatother
\crefname{originalreference}{original reference}{original references}
\Crefname{originalreference}{Original reference}{Original references}
\makeatletter
\newlabel{spaces-properties-lemma-points-cartesian}{{03H4}{source}{}{sourcepointer-9}{}}
\newlabel{spaces-properties-lemma-points-cartesian@cref}{{[originalreference][0][]03H4}{[1][0][]source}{}{}{}}
\makeatother
\title{Six equivalent characterizations of combined fibre finiteness over a point}
\author{The Stacks Project Authors}
\date{stacks-project:decent-spaces.tex:local-snapshot-20261012; GNU FDL 1.2 or later}
\begin{document}
\maketitle
\noindent\textit{Editorial note.} The source marks (4) implies (5) 'immediate' and (6) implies (5) 'clear'; both are single weakenings of the same statement form, carried as explicit obligations.\par\medskip
\noindent\textit{Allowed standard prerequisites (editorial).}\begin{itemize}
\item Weak characterization. lemma-weak-UR-finite-above-x (carded L26699). Gives (1) iff (3).
\item U-finiteness characterization. lemma-U-finite-above-x (carded L26697). Used in (2) implies (1).
\item R-finiteness characterization. lemma-R-finite-above-x (carded L26698). Used in (2) implies (1).
\item Quasi-compactness is etale-local on the target. Morphisms of Spaces, lemma quasi-compact-local (cross-chapter, named). Gives (3) implies (4).
\item Factorization of field points in the same class. Properties of Spaces, lemma equivalence-class-point-monomorphism (cross-chapter, named). Gives (4) implies (6).
\end{itemize}\medskip
\ifdefined\hypertarget\hypertarget{sourcepointer-1}{}\fi
\noindent\textit{Source reference.} Original lemma 03JU: sources/stacks-decent-spaces/pinned-original-source.tex, source lines 534--534. Navigation pointer only: the counted proof cites this statement by \textbackslash{}ref as a standing input; the deduction the counted proof has to make is carried out inside the counted span, and no part of that deduction is delegated to this pointer.\par\medskip
\ifdefined\hypertarget\hypertarget{sourcepointer-2}{}\fi
\noindent\textit{Source reference.} Original lemma 03JS: sources/stacks-decent-spaces/pinned-original-source.tex, source lines 485--485. Navigation pointer only: the counted proof cites this statement by \textbackslash{}ref as a standing input; the deduction the counted proof has to make is carried out inside the counted span, and no part of that deduction is delegated to this pointer.\par\medskip
\ifdefined\hypertarget\hypertarget{sourcepointer-3}{}\fi
\noindent\textit{Source reference.} Original lemma 040U: sources/stacks-decent-spaces/pinned-original-source.tex, source lines 687--687. Navigation pointer only: the counted proof cites this statement by \textbackslash{}ref as a standing input; the deduction the counted proof has to make is carried out inside the counted span, and no part of that deduction is delegated to this pointer.\par\medskip
\ifdefined\hypertarget\hypertarget{sourcepointer-4}{}\fi
\noindent\textit{Source reference.} Original remark 03II: sources/stacks-decent-spaces/pinned-original-source.tex, source lines 459--459. Navigation pointer only: the counted proof cites this statement by \textbackslash{}ref as a standing input; the deduction the counted proof has to make is carried out inside the counted span, and no part of that deduction is delegated to this pointer.\par\medskip
\ifdefined\hypertarget\hypertarget{sourcepointer-5}{}\fi
\noindent\textit{Source reference.} Original reference 03DP: focused-seeds/geometry/S9-decent-ur-finite-above-x/ref-tags.txt, source lines 5--5. Navigation pointer only: the counted proof cites this statement by \textbackslash{}ref as a standing input; the deduction the counted proof has to make is carried out inside the counted span, and no part of that deduction is delegated to this pointer.\par\medskip
\ifdefined\hypertarget\hypertarget{sourcepointer-6}{}\fi
\noindent\textit{Source reference.} Original reference 03HG: focused-seeds/geometry/S9-decent-ur-finite-above-x/ref-tags.txt, source lines 6--6. Navigation pointer only: the counted proof cites this statement by \textbackslash{}ref as a standing input; the deduction the counted proof has to make is carried out inside the counted span, and no part of that deduction is delegated to this pointer.\par\medskip
\ifdefined\hypertarget\hypertarget{sourcepointer-7}{}\fi
\noindent\textit{Source reference.} Original reference 03KG: focused-seeds/geometry/S9-decent-ur-finite-above-x/ref-tags.txt, source lines 7--7. Navigation pointer only: the counted proof cites this statement by \textbackslash{}ref as a standing input; the deduction the counted proof has to make is carried out inside the counted span, and no part of that deduction is delegated to this pointer.\par\medskip
\ifdefined\hypertarget\hypertarget{sourcepointer-8}{}\fi
\noindent\textit{Source reference.} Original reference 0H2X: focused-seeds/geometry/S9-decent-ur-finite-above-x/ref-tags.txt, source lines 8--8. Navigation pointer only: the counted proof cites this statement by \textbackslash{}ref as a standing input; the deduction the counted proof has to make is carried out inside the counted span, and no part of that deduction is delegated to this pointer.\par\medskip
\ifdefined\hypertarget\hypertarget{sourcepointer-9}{}\fi
\noindent\textit{Source reference.} Original reference 03H4: focused-seeds/geometry/S9-decent-ur-finite-above-x/ref-tags.txt, source lines 9--9. Navigation pointer only: the counted proof cites this statement by \textbackslash{}ref as a standing input; the deduction the counted proof has to make is carried out inside the counted span, and no part of that deduction is delegated to this pointer.\par\medskip
\begin{lemma}
\label{lemma-UR-finite-above-x}
Let $S$ be a scheme. Let $X$ be an algebraic space over $S$.
Let $x \in |X|$. The following are equivalent:
\begin{enumerate}
\item for every affine scheme $U$, any \'etale morphism
$\varphi : U \to X$ setting $R = U \times_X U$ the fibres of both
$$
|U| \longrightarrow |X|
\quad\text{and}\quad
|R| \longrightarrow |X|
$$
over $x$ are finite,
\item there exist schemes $U_i$ and \'etale morphisms
$U_i \to X$ such that $\coprod U_i \to X$ is surjective and for each
$i$, setting $R_i = U_i \times_X U_i$ the fibres of both
$$
|U_i| \longrightarrow |X|
\quad\text{and}\quad
|R_i| \longrightarrow |X|
$$
over $x$ are finite,
\item there exists a monomorphism $\Spec(k) \to X$ with $k$ a field
in the equivalence class of $x$, and for any affine scheme $U$ and \'etale
morphism $U \to X$ the fibre product $\Spec(k) \times_X U$ is
a finite scheme over $k$,
\item there exists a quasi-compact monomorphism $\Spec(k) \to X$
with $k$ a field in the equivalence class of $x$,
\item there exists a quasi-compact morphism $\Spec(k) \to X$
with $k$ a field in the equivalence class of $x$, and
\item every morphism $\Spec(k) \to X$ with $k$ a field in the
equivalence class of $x$ is quasi-compact.
\end{enumerate}
\end{lemma}

\begin{proof}
The equivalence of (1) and (3) follows on applying
Lemma \ref{lemma-weak-UR-finite-above-x}
to every \'etale morphism $U \to X$ with $U$ affine.
It is clear that (3) implies (2).
Assume $U_i \to X$ and $R_i$ are as in (2). We conclude from
Lemma \ref{lemma-U-finite-above-x}
that for any affine scheme $U$ and \'etale morphism $U \to X$
the fibre of $|U| \to |X|$ over $x$ is finite.
Say this fibre is $\{u_1, \ldots, u_n\}$. Then, as
Lemma \ref{lemma-R-finite-above-x} (1)
applies to $U_i \to X$ for some $i$ such that $x$ is in the image of
$|U_i| \to |X|$, we see that the fibre of
$|R = U \times_X U| \to |U| \times_{|X|} |U|$
is finite over $(u_a, u_b)$, $a, b \in \{1, \ldots, n\}$.
Hence the fibre of $|R| \to |X|$ over $x$ is finite.
In this way we see that (1) holds. At this point we know that
(1), (2), and (3) are equivalent.

\medskip\noindent
If (4) holds, then for any affine scheme $U$ and \'etale morphism
$U \to X$ the scheme $\Spec(k) \times_X U$ is on the one hand
\'etale over $k$ (hence a disjoint union of spectra of finite separable
extensions of $k$ by
Remark \ref{remark-recall})
and on the other hand quasi-compact over $U$ (hence quasi-compact).
Thus we see that (3) holds.
Conversely, if $U_i \to X$ is as in (2) and $\Spec(k) \to X$
is a monomorphism as in (3), then
$$
\coprod \Spec(k) \times_X U_i
\longrightarrow
\coprod U_i
$$
is quasi-compact (because over each $U_i$ we see that
$\Spec(k) \times_X U_i$ is a finite disjoint union spectra of fields).
Thus $\Spec(k) \to X$ is quasi-compact by
Morphisms of Spaces, Lemma \ref{spaces-morphisms-lemma-quasi-compact-local}.

\medskip\noindent
It is immediate that (4) implies (5). Conversely, let $\Spec(k) \to X$
be a quasi-compact morphism in the equivalence class of $x$. Let $U \to X$
be an \'etale morphism with $U$ affine. Consider the fibre product
$$
\xymatrix{
F \ar[r] \ar[d] & U \ar[d] \\
\Spec(k) \ar[r] & X
}
$$
Then $F \to U$ is quasi-compact, hence $F$ is quasi-compact.
On the other hand, $F \to \Spec(k)$ is \'etale, hence $F$ is a
finite disjoint union of spectra of finite separable extensions of $k$
(Remark \ref{remark-recall}). Since the image of $|F| \to |U|$
is the fibre of $|U| \to |X|$ over $x$ (Properties of Spaces, Lemma
\ref{spaces-properties-lemma-points-cartesian}), we conclude that
the fibre of $|U| \to |X|$ over $x$ is finite. The scheme
$F \times_{\Spec(k)} F$ is also a finite union of spectra of fields
because it is also quasi-compact and \'etale over $\Spec(k)$.
There is a monomorphism
$F \times_X F \to F \times_{\Spec(k)} F$, hence $F \times_X F$ is
a finite disjoint union of spectra of fields
(Schemes, Lemma \ref{schemes-lemma-mono-towards-spec-field}).
Thus the image of $F \times_X F \to U \times_X U = R$ is finite.
Since this image is the fibre of $|R| \to |X|$ over $x$ by
Properties of Spaces, Lemma \ref{spaces-properties-lemma-points-cartesian}
we conclude that (1) holds. At this point we know that
(1) -- (5) are equivalent.

\medskip\noindent
It is clear that (6) implies (5). Conversely, assume $\Spec(k) \to X$ is
as in (4) and let $\Spec(k') \to X$ be another morphism with $k'$ a
field in the equivalence class of $x$. By
Properties of Spaces, Lemma
\ref{spaces-properties-lemma-equivalence-class-point-monomorphism}
we have a factorization $\Spec(k') \to \Spec(k) \to X$ of the
given morphism. This is a composition of quasi-compact
morphisms and hence quasi-compact (Morphisms of Spaces,
Lemma \ref{spaces-morphisms-lemma-composition-quasi-compact}) as desired.
\end{proof}

\end{document}
